\PrelimChapter{Glosario}{\topBanner}\label{cap:glosario}
\vspace*{0.7cm}

En este apartado se encuentran términos clave y conceptos relevantes utilizados a lo largo de este proyecto.

\section*{A}
\begin{description}
	\item[Alta disponibilidad:] Diseño de un sistema orientado a minimizar los tiempos de inactividad, de manera que el servicio permanezca operativo de forma continua incluso ante la falla de alguno de sus componentes, generalmente mediante esquemas de redundancia \citep{IBM2026}.
	\item[Autenticación:] Proceso mediante el cual se verifica la identidad reclamada por un usuario o sistema, comprobando que los autenticadores presentados sean válidos \citep{NIST80063v4}.
	\item[Autenticación federada:] Mecanismo que permite a los usuarios acceder a recursos de múltiples sistemas o instituciones utilizando las credenciales de su organización de origen, sin necesidad de gestionar identidades duplicadas \citep{Prout2019}.
	\item[Autorización:] Proceso que determina si una acción o servicio solicitado está aprobado para una entidad específica, definiendo a qué recursos puede acceder y qué operaciones puede realizar sobre ellos \citep{OWASP2026}.
\end{description}

\section*{B}
\begin{description}
	\item[Balanceador de carga:] Componente que distribuye las solicitudes entrantes entre varios servidores con el fin de optimizar el uso de los recursos, evitar la sobrecarga de un nodo y garantizar la disponibilidad del servicio \citep{Cloudflare2026}.
\end{description}

\section*{C}
\begin{description}
	\item[Credencial:] Información que se presenta para verificar la identidad reclamada durante un proceso de autenticación, como contraseñas, certificados digitales o tokens \citep{NIST80063v4}.
\end{description}

\section*{M}
\begin{description}
	\item[Mapeo sistemático:] Metodología de revisión de la literatura que permite identificar, clasificar y analizar los estudios relevantes publicados sobre un área de investigación, ofreciendo una visión general del estado del arte \citep{Petersen2015}.
\end{description}

\section*{P}
\begin{description}
	\item[Principio de mínimo privilegio:] Práctica de seguridad que consiste en otorgar a cada usuario o proceso únicamente los permisos indispensables para ejecutar sus funciones autorizadas \citep{OWASP2026}.
\end{description}

\section*{S}
\begin{description}
	\item[Servicio de directorio:] Estructura jerárquica destinada al almacenamiento y consulta centralizada de información relacionada con objetos dentro de una red, como usuarios, grupos, recursos y políticas de acceso \citep{Microsoft2025}.
	\item[Software de código abierto:] Software cuyo código fuente se encuentra disponible para su consulta, modificación y redistribución bajo una licencia que cumple con la definición de código abierto, promoviendo la colaboración y la transparencia \citep{OSI2026}.
	\item[Software libre:] Software cuya licencia garantiza a sus usuarios las libertades de ejecutar, estudiar, modificar y distribuir el programa y sus versiones modificadas \citep{FSF2026}.
\end{description}

\section*{T}
\begin{description}
	\item[Tolerancia a fallos:] Capacidad de un sistema para continuar operando de manera continua después de que fallan uno o más de sus componentes críticos \citep{IBM2026}.
\end{description}
